(1) 固相中 $T_t=\alpha T_{xx}$，边界为 $T(0,t)=T_s$、$T(s(t),t)=T_m$，初始界面 $s(0)=0$。液相按模型保持 $T_m$，不输送热流；Stefan 潜热条件为 $\rho\mathcal L\,\dot s(t)=k T_x(s(t)^-,t)$，右侧为正的向冷壁导出的热流大小。
(2) 取相似变量 $\eta=x/(2\sqrt{\alpha t})$，热方程化为 $f''+2\eta f'=0$，其积分由误差函数给出。满足两处温度边界的解为 $T(x,t)=T_s+(T_m-T_s)\operatorname{erf}[x/(2\sqrt{\alpha t})]/\operatorname{erf}\lambda$，同时 $s(t)=2\lambda\sqrt{\alpha t}$。
(3) 界面内侧梯度为 $T_x(s^-,t)=(T_m-T_s)e^{-\lambda^2}/[\sqrt{\pi\alpha t}\operatorname{erf}\lambda]$，而 $\dot s=\lambda\sqrt{\alpha/t}$。代入 Stefan 条件及 $\alpha=k/(\rho c)$ 得 $\lambda e^{\lambda^2}\operatorname{erf}\lambda=c(T_m-T_s)/(\mathcal L\sqrt\pi)$；左侧对 $\lambda>0$ 严格增大，因而该正根唯一。
(4) 冷壁热通量模 $q_0(t)=k(T_m-T_s)/[\sqrt{\pi\alpha t}\operatorname{erf}\lambda]$，累计抽热 $Q_0/A=\int_0^tq_0(t')dt'=2\rho c(T_m-T_s)\sqrt{\alpha t}/(\sqrt\pi\operatorname{erf}\lambda)$。凝固潜热为 $Q_L/A=\rho\mathcal L\,s(t)=2\rho\mathcal L\lambda\sqrt{\alpha t}$；新固相显热降低为 $Q_s/A=\rho c\int_0^{s(t)}[T_m-T(x,t)]dx=2\rho c(T_m-T_s)\sqrt{\alpha t}(1-e^{-\lambda^2})/(\sqrt\pi\operatorname{erf}\lambda)$。借助界面方程，$Q_0=Q_L+Q_s$。
